\begin{abstract}
Se analizó experimentalmente el intercambio de energía térmica y el establecimiento del equilibrio mediante dos experiencias de calorimetría. En la primera se mezclaron porciones de agua a distintas temperaturas para estimar la capacidad calorífica del calorímetro; en la segunda se estudió la mezcla de agua y hielo, incluyendo la fusión. Se aplicaron balances basados en la conservación de energía y propagación conservadora de incertidumbres. Se obtuvo $C_{\mathrm{cal}}=(51.2\pm208.4)\,\mathrm{cal}/{}^{\circ}\mathrm{C}$ únicamente del tercer ensayo, pues en los dos primeros la temperatura final coincidió con la inicial del agua contenida en el calorímetro y anuló el denominador del cálculo. La temperatura experimental del sistema agua-hielo fue $(21\pm1)\,{}^{\circ}\mathrm{C}$, frente a predicciones de $20.42\,{}^{\circ}\mathrm{C}$ al considerar el calorímetro y $18.99\,{}^{\circ}\mathrm{C}$ al despreciarlo, con errores porcentuales de $2.86\%$ y $10.59\%$, respectivamente. La elevada incertidumbre de la calibración y la limitada precisión de las temperaturas restringieron la comparación. Aunque incluir el calorímetro redujo el error central y la fusión constituyó la principal contribución al calor absorbido por el hielo, los datos no permitieron validar con precisión el modelo de aislamiento térmico.
\end{abstract}

\begin{IEEEkeywords}
Calorimetría, equilibrio térmico, capacidad calorífica, conservación de energía, calor latente, incertidumbre experimental.
\end{IEEEkeywords}
